% Estilos compartidos por los diagramas de las Figs. 1 a 3.
\usetikzlibrary{positioning, arrows.meta, calc, fit, backgrounds}

\tikzset{
  diagram/.style={
    font=\scriptsize,
    >={Stealth[length=1.4mm, width=1.0mm]},
    line width=0.4pt,
    node distance=3mm,
  },
  block/.style={
    draw=black!65, fill=blue!8, rounded corners=1.2pt,
    align=center, inner sep=2pt, minimum height=9mm,
  },
  auxblock/.style={block, fill=black!6},
  lossblock/.style={block, fill=orange!14, minimum height=6mm},
  note/.style={align=center, font=\tiny, text=black!65, inner sep=1pt},
  flow/.style={->, draw=black!75},
  aux/.style={->, draw=black!55, densely dashed},
  % vineta con un fotograma real del material del articulo
  photo/.style={
    inner sep=0pt, outer sep=0pt,
    draw=black!65, line width=0.4pt,
  },
}

% ancho de las vinetas de los diagramas
\newcommand{\thumbw}{0.95cm}

% Fotograma suelto: \singleframe{nombre}{archivo}{posicion}
\newcommand{\singleframe}[3]{%
  \node[photo] (#1) at #3 {\includegraphics[width=\thumbw]{#2}};%
}

% Antes habia una pila de hojas detras de cada fotograma para sugerir que era
% una secuencia, pero a un centimetro de ancho se leia como recuadros sueltos
% con la foto encima. Un solo fotograma con borde se entiende mejor; la idea de
% secuencia la lleva el rotulo.
